\chapter{Goal, quadratic, and nonlinear programming}
\label{chap:optimizacion-avanzada}

This chapter brings together three distinct extensions of linear programming.
Goals order multiple objectives; quadratic programming adds curvature to the
objective while retaining linear constraints; nonlinear programming allows
general expressions. Choosing a family changes both the formulation and the
guarantee the solver can offer. Comparing the number displayed as the objective
is insufficient: priorities, feasibility, and exactness must also be reviewed.

\section{Goal programming}
\subsection{Priorities and deviations}
An organization may first seek to meet a goal and then optimize another
quantity without sacrificing what has already been achieved. In lexicographic
programming, \(g_1(x)\) is optimized, its optimal value is fixed, and
\(g_2(x)\) is optimized within that set; the process continues in the declared
order. Objectives are not automatically added with large weights. The order
expresses a policy decision: an improvement in a later priority never
compensates for deterioration in an earlier one.

For a goal \(a^\top x=t\), the deviation variables satisfy
\[
 a^\top x+n-p=t,\qquad n,p\geq0.
\]
\(n\) measures shortfall and \(p\) surplus. Minimizing \(n+p\) penalizes both
directions; minimizing only \(n\) tolerates surplus. Coefficients and units
must express the intended penalty. QSB does not infer these equations:
the user declares variables, objectives, and constraints.

\subsection{Example: protecting a production commitment}
There is capacity for ten units shared between products \(A\) and \(B\).
The first aim is to bring \(A\) close to eight units, followed by maximizing \(B\):
\[
\begin{aligned}
 &\text{Priority 1: }\min(n_1+p_1),\\
 &\text{Priority 2: }\max B,\\
 &A+B\leq10,\qquad A+n_1-p_1=8,\\
 &A,B,n_1,p_1\geq0.
\end{aligned}
\]
The priority does not itself require exact fulfillment if the goal is
impossible; it minimizes deviation subject to the hard constraints.
In this example, zero deviation is attainable. Once it is fixed,
\(A=8\), so the subsequent maximum is \(B=2\).

The file \path{manual/examples/opt_goal.json} contains all variables,
the two rows, and the ordered \texttt{goals} list; each priority has a name,
\texttt{sense}, and \texttt{coefficients}. Bounds and types are declared in
arrays aligned with \texttt{variableNames}. Run:
\begin{lstlisting}
swift run qsb validate-goal-json manual/examples/opt_goal.json
swift run qsb solve-goal-json manual/examples/opt_goal.json --backend native
\end{lstlisting}
The output document retains the model, \texttt{backend} metadata, and
solution.

\texttt{goalOutcomes} displays priority, sense, and value.
\begin{table}[htbp]
\centering
\caption{Lexicographic result for the production case}
\begin{tabular}{lr}
\toprule
Quantity & Value \\
\midrule
Priority 1: total deviation & 0 \\
Priority 2: production of \(B\) & 2 \\
\(A\) & 8 \\
\(n_1\), \(p_1\) & 0, 0 \\
\bottomrule
\end{tabular}
\end{table}
Capacity is active and the goal is met. Reversing the order makes the first
objective yield \(B=10\); with that value fixed, \(A=0\), \(n_1=8\), and
deviation is eight. This contrast invites discussion of who determines
precedence and which commitment it protects.

\subsection{macOS workflow and solver character}
Open the normalized file in \texttt{Goal Programming}. In
\texttt{Definition}, review variables and bounds, the priority list, and
linear constraints. The move-up and move-down controls change a priority's
position. Enter the deviation equation as an equality and verify the
negative sign of \(p_1\). Press \texttt{Validate} and \texttt{Run}; the native
priority view shows the ordered objectives, final quantities, and backend
information. JSON provides the complete exportable record. Tables support
copying and pasting; coefficient columns must match the current variable
order. Large tables use row selection and coefficient groups.

The native backend solves a sequence of LP or integer problems, adding an
equality to preserve each previous optimum. The continuous case is labeled
\texttt{exact}; this term characterizes the method, subject to floating-point
tolerances. Integer goals use an educational search for small problems,
identified as \texttt{fixtureScale}, and depend on branch-and-bound limitations.
The special relaxation of integer deviation pairs applies only when integral
equations justify integral residuals; it does not remove integrality from
the decisions.

\texttt{validate} does not solve. \texttt{external} executes each priority
with HiGHS if installed; QSB preserves the lexicographic order and fixing.
A missing executable causes an explicit error. An infeasible later priority,
an unbounded objective, or exhausted limits require reviewing the diagnostic
rather than interpreting partial results as the final plan.

\section{Quadratic programming}
\subsection{Objective convention and assumptions}
QSB represents the objective as
\[
 f(x)=x^\top Qx+c^\top x,
\]
with linear constraints, bounds, and variable types. JSON uses
\texttt{quadraticMatrix} for \(Q\) and \texttt{linearCoefficients} for \(c\).
This convention has no factor of \(1/2\). For a symmetric matrix, the
coefficient of the product \(x_ix_j\), \(i\ne j\), is \(2Q_{ij}\): a term
\(4XY\) is represented by \(Q_{12}=Q_{21}=2\). Do not place 4 in both
cells, since that would double the effect.

A symmetric positive-definite matrix produces a strictly convex objective;
its minimum over a convex region, if it exists, is unique. A negative-definite
matrix produces the equivalent strictly concave maximization case.
Stationarity and feasibility conditions then identify the global optimum.
An indefinite matrix changes the problem and is not covered by that guarantee.

\subsection{Example: adjustment with limited capacity}
Two quantities should be brought close to 2 and 3, respectively, with total
capacity four. Removing the constant 13 from the squared distance gives
\[
\begin{aligned}
 \min\quad & X^2+Y^2-4X-6Y\\
 \text{subject to}\quad &X+Y\leq4,\qquad X,Y\geq0.
\end{aligned}
\]
The matrix is the identity and the linear vector is \((-4,-6)\). Removing
the constant preserves optimal quantities but changes the reported objective
value: the original squared distance is \(f(X,Y)+13\).

The model is in \path{manual/examples/opt_qp.json}. Run:
\begin{lstlisting}
swift run qsb validate-qp-json manual/examples/opt_qp.json
swift run qsb solve-qp-json manual/examples/opt_qp.json --backend native
\end{lstlisting}
The optimum is \(X=1.5\), \(Y=2.5\), \(f=-12.5\), with
\texttt{Capacity} in \texttt{activeConstraints}. The original squared distance
is 0.5. A negative objective is perfectly valid: it results from the omitted
constant and does not represent a negative distance.

Without the capacity constraint the solution would be \((2,3)\), which is
infeasible. With multiplier \(\lambda\geq0\), the stationarity equations are
\[
 2X-4+\lambda=0,\qquad2Y-6+\lambda=0.
\]
The active capacity requires \(X+Y=4\), giving \(\lambda=1\), \(X=1.5\), and
\(Y=2.5\). The lower bounds are inactive. Strict convexity and the conditions
above certify the global minimum of the continuous model.
QSB displays active constraints, but this multiplier derivation is an
independent academic check.

\subsection{macOS workflow and limitations}
Open the JSON in \texttt{Quadratic Programming}. \texttt{Definition} separates
variables and bounds, the quadratic matrix, and linear constraints. Enter
the linear coefficients in the variable rows and the identity in the matrix;
confirm the names on both axes. Add \texttt{Capacity} with coefficients
\((1,1)\), relation \(\leq\), and right-hand side 4. Validate and press
\texttt{Run}. The native solution displays the objective, variables, bounds,
active constraints, and exactness; JSON allows inspection of the document.

\texttt{Copy Table} and \texttt{Paste Table} transfer the complete matrix.
Reordering known variables preserves both axes by name; pasted headers must
match the current order. For more than twelve variables, the editor displays
one row and groups of columns, although table exchange remains complete.
This editor capability does not imply support for large problems in the
native solver.

The continuous native solver enumerates active sets and solves KKT systems
for strictly convex minimization or strictly concave maximization. Its label
is \texttt{exact}, with numerical calculations. The search rejects more than
twenty candidate inequalities, counting bounds that it converts into
inequalities; the cost of set enumeration grows rapidly. The integer native
solver supports only diagonal concave maximization with nonnegative variables,
up to four variables, and enumerates points within derived finite bounds.
It is identified as \texttt{fixtureScale}; it is not a general MIQP solver.

The external HiGHS backend supports translation of the supported strictly
convex/concave continuous cases. QSB adapts the convention using a Hessian
\(H=2Q\), because the external solver uses \(\tfrac12 x^\top Hx\).
The adaptation also preserves sense, bounds, and constraints, and verifies
the objective and feasibility when reading the solution.
HiGHS does not solve MIQP in this integration; quadratic integer variables
require the native scope described above or a future engine. The
\texttt{validate} mode does not compute an optimum.

\section{Nonlinear programming}
\subsection{Expressions and domain}
Nonlinear programming (NLP) allows an objective \(f(x)\) and constraints
\(g_i(x)\ \{\leq,=,\geq\}\ b_i\). QSB accepts textual expressions,
ordered variables, and nullable bounds. The evaluator recognizes addition,
subtraction, multiplication, division, powers, parentheses, implicit
multiplication, and \texttt{exp}, \texttt{log}, \texttt{sqrt}, \texttt{sin},
\texttt{cos}. Use explicit multiplication in teaching materials to aid reading.
Names are recognized without case sensitivity; consistent spelling is advisable.

Valid syntax does not guarantee a valid domain at every point:
\texttt{log(X)} requires \(X>0\), \texttt{sqrt(X)} requires \(X\geq0\), and
division requires a nonzero denominator. Bounds must correspond to the case
and protect the domain where possible. Sine and cosine angles are expressed
in radians. JSON stores \texttt{objectiveExpression} and, for each constraint,
\texttt{expression}, \texttt{relation}, and \texttt{rhs}.

\subsection{Example: constrained distance}
The file \path{manual/examples/opt_nlp.json} specifies
\[
 \min\ (X-3)^2+(Y-2)^2,\qquad X+Y\leq4,\qquad0\leq X,Y\leq4.
\]
The objective is the squared distance to point \((3,2)\), which violates
capacity. Its projection onto \(X+Y=4\) is \((2.5,1.5)\), and the mathematical
value is 0.5. Although this case could be formulated as QP, it serves to
compare expression input and an approximate method with a known analytical
solution.
\begin{lstlisting}
swift run qsb validate-nlp-json manual/examples/opt_nlp.json
swift run qsb solve-nlp-json manual/examples/opt_nlp.json --backend native
\end{lstlisting}
The verified execution returns approximately
\(X=2.500000002\), \(Y=1.499999999\), objective 0.4999999995, and maximum
violation \(5.0\times10^{-10}\). The evaluation of \texttt{Capacity} is
approximately 4.0000000005. The \texttt{iterations} field records 361 cumulative
iterations of the selected search; it does not count all evaluations of
searches from the other initial points. This value describes this execution,
not a universal limit or a measure comparable with simplex iterations.

The small violation does not authorize physically exceeding capacity. It is
a numerical search residual; for real implementation, check the values
against the tolerance and units of the case. The output contains
\texttt{constraintEvaluations}, \texttt{maximumViolation},
\texttt{iterations}, the objective, and variable values, as well as the model
and metadata. Reporting the objective without the residual would hide
relevant information.

\subsection{macOS use and approximate guarantee}
Open the JSON in \texttt{Nonlinear Programming}. The definition allows editing
the sense, objective expression, bounded variables, and constraints with an
expression, relation, and right-hand side. Enter
\texttt{(X-3)\^{}2+(Y-2)\^{}2}, bounds 0 and 4, and \texttt{X+Y} with
\(\leq4\). Use \texttt{Validate} before \texttt{Run}. An invalid expression is
retained for correction. Variables and constraints support spreadsheet
exchange; each valid paste replaces a complete table.
The solution displays the candidate, bounds, evaluations, and maximum
violation, together with its approximate character and exportable JSON.

The native backend uses automatic differentiation, searches from multiple
points, and progressive constraint penalties. It is deterministic, is limited
to four variables, and is identified as \texttt{approximate}. It also includes
a low-dimensional search on equality manifolds for small cases.
These mechanisms do not prove global optimality for nonconvex functions.
A feasible candidate may be local, and a failed search does not in itself
constitute a mathematical proof of infeasibility.

Strict inequalities from imports are normalized to non-strict boundaries
and generate a structured warning. The
\texttt{normalizedStrictInequalities} field records that transformation;
it must be reviewed if the boundary has a physical meaning excluded by the
case. \texttt{validate} diagnoses without solving. \texttt{external} is
explicitly unavailable for NLP; select the native backend for supported
educational problems.

\section{How to compare and review results}
Retain each input together with its output and the backend used. For goals,
review the order and each fixed value. For QP, check the matrix convention
and supported curvature. For NLP, review domains, bounds, evaluations, and
maximum violation before describing a solution as implementable. File
validation and search convergence provide different kinds of evidence.

All three families have a native definition view and a JSON view; unfinished
fields remain editable when validation fails. In native macOS runs, solving
takes place in the background and duplicate runs are prevented while one is
in progress. Changing the model can make that calculation's result stale.
There is currently no cancellation control for these workers. Accessible
labels on controls and results do not replace a complete VoiceOver and
keyboard navigation audit, which remains pending for these editors.

\section{Exercises}
\begin{enumerate}
 \item Reverse the priorities in the goal example. Verify \(B=10\),
 \(A=0\), and total deviation 8. Explain which commitment has changed.
 \item Reduce the capacity in the goal example to 6. Calculate the minimum
 deviation of the first priority and justify why the second cannot recover
 production of \(B\) without worsening the first.
 \item In the QP, increase capacity to 5. Predict \((X,Y)=(2,3)\),
 objective \(-13\), and zero original distance. Check against QSB.
 \item Introduce a term \(4XY\) in a test objective. Write the correct
 symmetric matrix and examine its curvature before choosing a solver.
 \item Reformulate the example NLP objective in QP form, retaining
 constraints and bounds. Compare quantities and residuals; explain the
 difference between its objective with a constant and the quadratic
 objective without it.
 \item Propose a one-variable function with two local minima. Explain what
 additional evidence you would need to attribute global optimality to a
 candidate obtained by the native NLP solver.
\end{enumerate}
